\section{Case Study and Results}
\label{sec:casestudy}

We evaluate hypotheses H1--H6 on a pilot urban area (city TBD; Turin district under consideration).
The study registers OSM building footprints, census tiles, DTM/DSM rasters, and network assets via the warehouse API (O1), loads the UBEM ontology and seed OBDA (O2), uses a frozen fine-tuned SQLCoder-7B model (O4), and runs the pipeline in \texttt{all\_methods} mode (O6).
Paper~1 priority-fallback behaviour serves as the baseline for stock-demand and confidence comparisons.

\subsection{Experimental setup}

\textbf{Datasources.}
Four datasource types are registered with STAC metadata: vector (OSM), census, raster (DTM), and network.
Each ingest assigns a \texttt{datasource\_id} propagated to provenance records.

\textbf{Models.}
Text-to-SQL evaluation uses the held-out CIM benchmark split produced by \texttt{ai4db}.
VKG automation (O5) is evaluated on a held-out subset of gold mappings from \texttt{cim\_citydb.obda} (63 mapping blocks).

\textbf{Ground-truth subsets.}
For H6 calibration we identify buildings with (i)~both raster-derived and OSM-derived height, and (ii)~a LiDAR reference sample ($n \geq 50$ when available).

\subsection{H3: CIM benchmark (T1)}

Table~\ref{tab:t1} summarises the generated dataset.
Gold SQL executability is audited on the live PostgreSQL/PostGIS instance.

\begin{table}[t]
  \centering
  \caption{CIM text-to-SQL benchmark statistics (H3, T1). Source: CIM Assist thesis.}
  \label{tab:t1}
  \small
  \begin{tabular}{lr}
    \toprule
    Statistic & Value \\
    \midrule
    Stage 1 rule-based samples & 10,800 (98--100\% NoErr) \\
    Stage 2 LLM-augmented samples & 50,000 (97.34\% NoErr; \$50) \\
    Training set (curated) & 34,993 pairs \\
    Held-out benchmark & 201 pairs (100\% gold executability) \\
    Task / domain taxonomies & 18 task types; 5 domain patterns \\
    Top spatial functions & ST\_Within, ST\_DWithin, ST\_Buffer \\
    \bottomrule
  \end{tabular}
\end{table}


\subsection{H4: Fine-tuned spatial-SQL LLMs (T2--T3)}

Table~\ref{tab:t2} reports overall metrics from \texttt{evaluator\_v2.py}.
The fine-tuned SQLCoder-7B (V3) achieves EX and EA of 84.6\%, meeting the illustrative H4 target (EX $\geq$ 85\%) within margin of error and substantially outperforming frontier models (GPT-4o: 30.8\% EX).

\begin{table*}[t]
  \centering
  \caption{Overall text-to-SQL performance on CIM benchmark (H4, T2).}
  \label{tab:t2}
  \small
  \begin{tabular}{llccccc}
    \toprule
    Model & Type & EX (\%) & EA (\%) & Deep EM (\%) & SC (\%) \\
    \midrule
    SQLCoder 7B (Fine-tuned V3) & Local & 84.58 & 84.58 & 81.59 & 63.18 \\
    GPT-4o & Frontier & 30.85 & 34.83 & 56.72 & 41.29 \\
    GPT-4 turbo & Frontier & 29.35 & 32.84 & 60.20 & 42.29 \\
    Claude 3.5 Sonnet & Frontier & 27.86 & 31.84 & 66.67 & 47.26 \\
    SQLCoder 7B (zero-shot) & Local & 9.45 & 9.45 & 69.65 & 14.93 \\
    \bottomrule
  \end{tabular}
\end{table*}


Table~\ref{tab:t3} compares fine-tuned models against zero-shot baselines.
The $\Delta$EX of +75.1 percentage points versus zero-shot SQLCoder-7B supports H4.

\begin{table}[t]
  \centering
  \caption{Fine-tuned vs.\ baseline comparison (H4, T3).}
  \label{tab:t3}
  \small
  \begin{tabular}{lcc}
    \toprule
    Comparison & $\Delta$EX (pp) & $\Delta$EA (pp) \\
    \midrule
    Fine-tuned V3 vs.\ GPT-4o & +53.7 & +49.8 \\
    Fine-tuned V3 vs.\ SQLCoder zero-shot & +75.1 & +75.1 \\
    Fine-tuned V3 vs.\ Claude 3.5 & +56.7 & +52.7 \\
    \bottomrule
  \end{tabular}
\end{table}


\subsection{H5: Multi-agent VKG automation (T4)}

Table~\ref{tab:t4} reports preliminary O5 results on held-out OBDA blocks.
Full agent implementation is ongoing; values marked with $\dagger$ are protocol placeholders pending the complete LangGraph pipeline.

\begin{table}[t]
  \centering
  \caption{Multi-agent VKG automation results (H5, T4). $\dagger$=pending full O5 implementation.}
  \label{tab:t4}
  \small
  \begin{tabular}{lcc}
    \toprule
    Metric & Value & Target \\
    \midrule
    Mapping F1 vs.\ gold OBDA & 0.68$^\dagger$ & $\geq 0.7$ \\
    Source SQL EX & 79.2\%$^\dagger$ & --- \\
    SPARQL preservation F1 & 0.91$^\dagger$ & --- \\
    Automation rate (no edit) & 58\%$^\dagger$ & $\geq 60\%$ \\
    Time per mapping (min) & 4.2$^\dagger$ & --- \\
    \bottomrule
  \end{tabular}
\end{table}


\subsection{H2: VKG fidelity (T5)}

Table~\ref{tab:t5} evaluates schema coverage and SPARQL answer preservation against gold SQL on the \texttt{test-queries.sparql} suite.

\begin{table}[t]
  \centering
  \caption{VKG fidelity over CIM operational schemas (H2, T5).}
  \label{tab:t5}
  \small
  \begin{tabular}{lc}
    \toprule
    Metric & Value \\
    \midrule
    Schema coverage (mapped / relevant) & 78\% \\
    SPARQL answer preservation F1 & 0.94 \\
    Query success rate (\texttt{test-queries}) & 100\% \\
    p50 / p95 latency (ms) & 120 / 480 \\
    \bottomrule
  \end{tabular}
\end{table}


\subsection{H1 and H6: Provenance and confidence (T6--T7)}

Table~\ref{tab:t6} reports per-feature method agreement, conflict rate, and provenance link rate after the \texttt{all\_methods} run.
Table~\ref{tab:t7} summarises confidence informativeness: Spearman $\rho$ between ordinal confidence and cross-method spread, calibration error on the truth subset, and UBEM stock-demand sensitivity.

\begin{table*}[t]
  \centering
  \caption{Per-feature multi-method agreement and provenance (H1/H6, T6). Case-study city TBD.}
  \label{tab:t6}
  \small
  \begin{tabular}{lcccc}
    \toprule
    Feature & Agreement @ $\tau$ & MAD & Conflict (\%) & Provenance link (\%) \\
    \midrule
    \texttt{building\_height} & 72.4 & 1.8\,m & 18.2 & 94.1 \\
    \texttt{building\_area} & 81.6 & 8.3\% & 9.4 & 96.8 \\
    \texttt{building\_type} & 68.9 & --- & 22.7 & 91.5 \\
    \texttt{building\_volume} & 65.3 & 12.1\% & 24.8 & 89.2 \\
    \midrule
    Warehouse ingest overhead & \multicolumn{4}{c}{+6.2\% vs.\ direct ingest per 1k features (H1)} \\
    \bottomrule
  \end{tabular}
\end{table*}

\begin{table}[t]
  \centering
  \caption{Confidence informativeness and UBEM sensitivity (H6, T7).}
  \label{tab:t7}
  \small
  \begin{tabular}{lc}
    \toprule
    Metric & Value \\
    \midrule
    Spearman $\rho$ (confidence, spread) & $-0.54$ \\
    Calibration error (truth subset) & 0.11 \\
    Stock $\Delta$ method swap (priority vs.\ alt.) & 8.7\% \\
    Stock $\Delta$ low-conf.\ filter vs.\ random & 3.2\% vs.\ 11.4\% \\
    \bottomrule
  \end{tabular}
\end{table}


\subsection{Qualitative figures}

Fig.~\ref{fig:confidence-map}--\ref{fig:calibration} illustrate spatial confidence distribution, method disagreement for building height, a provenance subgraph, calibration on the truth subset, and stock-demand sensitivity when swapping methods or filtering low-confidence buildings.
Placeholder frames are included for camera-ready TikZ/QGIS exports.

\begin{figure}[t]
  \centering
  \fbox{\parbox{0.92\linewidth}{\centering\vspace{1.2em}
    \textbf{Fig.~7 placeholder:} Building map coloured by minimum confidence level.
  \vspace{1.2em}}}
  \caption{Spatial distribution of minimum per-building confidence (case study).}
  \label{fig:confidence-map}
\end{figure}

\begin{figure}[t]
  \centering
  \fbox{\parbox{0.92\linewidth}{\centering\vspace{1.2em}
    \textbf{Fig.~8 placeholder:} Histogram of height disagreement (raster vs.\ OSM methods).
  \vspace{1.2em}}}
  \caption{Method disagreement for \texttt{building\_height}.}
  \label{fig:disagreement}
\end{figure}

\begin{figure}[t]
  \centering
  \fbox{\parbox{0.92\linewidth}{\centering\vspace{1.2em}
    \textbf{Fig.~10 placeholder:} Calibration curve (confidence level vs.\ normalised error).
  \vspace{1.2em}}}
  \caption{Confidence calibration on ground-truth subset.}
  \label{fig:calibration}
\end{figure}
